\chapter*{Appendix} %{Appendices}
\label{appendices}
%DELETEME: everything that does not fit into your work, e.g. a 5 page table that breaks the reading flow, should be placed here

\todo{check bibstyle - check all relevant docs in Lit review folder are included}

\addcontentsline{toc}{section}{Product Branding Etymology}
%excursus
\section*{On Etymology of Product Names}
\label{etymology}

Product naming is a branding problem which companies deal with differently. In the case of Apple, with exception of the `Apple Watch' and on their consumer products line, they opt for a clever version using the letter `i' in front of simple common words, like iBook, iPad, iPhone. though this naming convention originally arose from `i' to stand for `interactive' celebrating the early adoption of the internet on the Macintosh product line, it has turned into a naming breakthrough since it had two effects: 1) it gives the user of Apple software or hardware product a feeling of personalisation for the letter `i' would make something sound as belonging to oneself, namely `` `I', as a person'' as a prefix for the name of the object create the name of the product, and 2) it is a marketing strategy to associate all products and services carrying the prefix `i' with Apple without the need to name the company's name. This applies to products and services that came to life even long after people's adoption of the internet and the widespread use of Apple products, like with iCloud as a web service. Another example is with professional software products like AutoCAD combining the company's name (AutoDesk) with the product's functionality  (Computer Aided Design). Companies like Adobe clearly prefer to name their products and services with their brand's full name included, as the case with Adobe Photoshop, Adobe Flash (formely know as MacroMedia Flash). Similarly, as the youngest of them, Amazon adopts a combined approach by stating the full name of the company on its consumer line like with `Amazon Prime', `Amazon Alexa' and a mix between `AWS' and `Amazon' like with `Amazon SageMaker', `AWS Firewall Manager' etc. on its developer B2B line.  


%\todo{did something go missing here? check git!}
%amazon transcri ?

\addcontentsline{toc}{section}{Lex vs. Alexa Skills}
\section*{Difference Between Lex and Alexa Skills}
\label{lexAlexa}  

%\todo{cite? indent? copied text}
Amazon Lex is a service for building conversational interfaces using voice and text. Powered by the same conversational engine as Alexa, Amazon Lex provides high quality speech recognition and language understanding capabilities, enabling addition of sophisticated, natural language \'chatbots\' to new and existing applications. Amazon Lex reduces multi-platform development effort, allowing you to easily publish your speech or text chatbots to mobile devices and multiple chat services, like Facebook Messenger, Slack, Kik, or Twilio SMS. Native interoperability with AWS Lambda, AWS MobileHub and Amazon CloudWatch and easy integration with many other services on the AWS platform including Amazon Cognito, and Amazon DynamoDB makes bot development effortless. \footnote{
More information: \url{https://stackoverflow.com/questions/42982159/differences-between-using-lex-and-alexa\#URL}\\
\url{https://aws.amazon.com/about-aws/whats-new/2017/09/export-your-amazon-lex-chatbot-to-the-alexa-skills-kit/}\\
\url{https://aws.amazon.com/lex/faqs/}
}









%%%%%%%%%%%%%%%%%%%%%%%%%%%%%%%%%%%%%%%%%%%%%%%%%%%%%%%%%%%%%%%%%%%%%%%%%%%%%
%%%%%%%%%%%%%%%%%%%%%%%%%%%%%%%%%%%%%%%%%%%%%%%%%%%%%%%%%%%%%%%%%%%%%%%%%%%%%
\clearpage

\section*{Documentation (Base) URLs}

\todo{use a base url to avoid extensive text in footnotes with href such that a click would still redirect right\\ -kbase \\ -Alexa blogs\\
\\
AWS Glossary\\ \t{a\t{ws}}\href{www.apple.com}{\lstinline|/doing-something\#hej|}
\url{https://docs.aws.amazon.com/general/latest/gr/glos-chap.html}

}



\subsection*{Amazon Web Services - \t{a\t{ws}}}


\begin{flushleft}
	\begin{tabularx}{\textwidth}{lX}




AWS Glossary &
\url{https://docs.aws.amazon.com/general/latest/gr/glos-chap.html}\\

\end{tabularx}
\end{flushleft}
	


\subsection*{Alexa Skills Kit - \t{a\t{sk}}}
\begin{flushleft}
	\begin{tabularx}{\textwidth}{l   X}

Documentation &
\url{https://developer.amazon.com/docs/custom-skills/} \\


Handling Requests &
\href{https://developer.amazon.com/docs/custom-skills/handle-requests-sent-by-alexa.html}{\t{a\t{sk}}/handle-requests-sent-by-alexa.html}

\end{tabularx}
\end{flushleft}

\subsection*{GitHub Repositories - \t{GH}}


\begin{flushleft}
	\begin{tabularx}{\textwidth}{lX}
		
		
A Skill Template with API Calls: & 
\url{https://github.com/skilltemplates/api-starter-alexa}\\

Node.js code &
\url{https://github.com/alexa/alexa-skills-kit-sdk-for-nodejs}\\

Building Responses in Node.js & 
\url{https://github.com/alexa/alexa-skills-kit-sdk-for-nodejs/wiki/Response-Building}\\


Alexa Cookbook (maintained by Amazon Team) & 
\url{https://github.com/alexa/alexa-cookbook/}\\

Handling Responses (Multi-turn conversation) & 
\url{https://github.com/alexa/alexa-cookbook/tree/master/handling-responses}\\
		
SSML Audio & 
\url{https://github.com/alexa/alexa-cookbook/tree/0afbfcb8ddcd911e83216b4528b20f3a1628c99f/handling-responses/ssml-audio}\\
		
	\end{tabularx}
\end{flushleft}



%\t{a\t{sk}}\href{   }{\lstinline|   |}

%find \url{https://developer.amazon.com/docs/custom-skills
% replace with \t{a\t{sk}}\href{\url{https://developer.amazon.com/docs/custom-skills}{\lstinline|   

%then |}


%%%%%%%%%%%%%%%%%%%%%%%%%%%%%%%%%%%%%%%%%%%%%%%%%%%%%%%%%%%%%%%%%%%%%%%%%%%%%
\clearpage

\addcontentsline{toc}{section}{JSON Query code structure}
\section*{Outtakes from a JSON query node}
\label{query:dl}

HTML Tags inside Text
\begin{minted}[tabsize=2, bgcolor=bgkolor, breaklines, fontsize=\footnotesize]{json}

{ "d115Description": "Eine Fiktionsbescheinigung wird ausgestellt, wenn über einen beantragten Aufenthaltstitel noch nicht entschieden werden kann, z. B. weil<br />\n<br />\n<ul class=\"list\"><li> ... Ausländerakte ... liegt,</li> ...
\end{minted}

Ambiguous words that can be a plus for the Virtual Citizen Assistant but become irrelevant or confusing for the Skill (e.g. \mintinline{text}{Fiktionsbescheinigung})
\begin{minted}[tabsize=2, bgcolor=bgkolor, breaklines, fontsize=\footnotesize]{json}
{... "d115Synonym": [
"Aufenthaltserlaubnis","Aufenthaltstitel","vorläufig",
"Visum","Fiktionsbescheinigung"] ...
\end{minted}

Inconsistent or unconventional text encoding (e.g. \mintinline{text}{::: false :::} )
\begin{minted}[tabsize=2, bgcolor=bgkolor, breaklines, fontsize=\footnotesize]{json}
{... "d115InfoLaw": ["§ 81 Aufenthaltsgesetz - AufenthG  ::: false ::: http://www.geset...net.de/aufenthg_2004/__81.html"] ...
\end{minted}

Text gets lost in translation from reading to speaking or is hard to read out loud
\begin{minted}[tabsize=2, bgcolor=bgkolor, breaklines, fontsize=\footnotesize]{json}
{... "d115Prerequisites": [ "Persönliche Vorsprache ist erforderlich ::: false ::: ",
"Rechtmäßiger Aufenthalt mit oder ohne Aufenthaltstitel ::: Zum Zeitpunkt des Antrags muss die Antragstellerin oder der Antragsteller entweder<br />\n<br />\n<ul class=\"list\"><li>einen gültigen Aufenthaltstitel (Aufenthaltserlaubnis oder nationales Visum für längerfristige Aufenthalte - Kategorie D - ) [...]",
"Antrag auf Aufenthaltstitel ::: Eine Fiktionsbescheinigung wird nur dann ausgestellt, wenn [...] ::: ",
"Hauptwohnsitz in Berlin ::: false ::: "] ...
\end{minted}

\begin{minted}[tabsize=2, bgcolor=bgkolor, breaklines, fontsize=\footnotesize]{json}
{... "d115Requirements": [ "Bisheriger Aufenthaltstitel ::: Soweit vorhanden, ist der bisherige Aufenthaltstitel mitzubringen, z.B. der elektronische Aufenthaltstitel (eAT). ::: ",
"Nachweis über den Hauptwohnsitz in Berlin ::: <ul class=\"list\"><li>Bescheinigung über die Anmeldung der Wohnung (Meldebestätigung) <strong>oder</strong></li><li>Mietvertrag und Einzugsbestätigung des Vermieters</li></ul>Mehr zum Thema im Abschnitt „Weiterführende Informationen“ ::: ",
"..."] ...
\end{minted}

Not a single readable item price
\begin{minted}[tabsize=2, bgcolor=bgkolor, breaklines, fontsize=\footnotesize]{json}
{..."d115Fees": "Ab dem 01.09.2017:<br />\n<br />\n<ul class=\"list\"><li>Für Erwachsene: 13,00 Euro</li><li>Für Minderjährige: 6,50 Euro</li></ul>Gebührenfrei für<br />\n<ul class=\"list\"><li>Türkische Staatsangehörige</li><li>Asylberechtigte</li><li>
Ausländer, die im Bundesgebiet die Rechtsstellung ausländischer Flüchtlinge oder subsidiär Schutzberechtigter genießen</li><li>Resettlement-
Flüchtlinge im Sinne von § 23 Abs. 4 S. 1 AufenthG</li></ul>" ...
\end{minted}

\begin{minted}[tabsize=2, bgcolor=bgkolor, breaklines, fontsize=\footnotesize]{json}
{..."d115ProcessTime": "Die Fiktionsbescheinigung wird bei Vorsprache ausgestellt." ...
\end{minted}

%
%"d115AppointmentLink": "",
%"d115ServiceLocations": [
%"121885",
%"327437",
%"327805"
%],
%"d115ServiceLocationsJson": [
%"{\"location\":\"327437\",\"url\":\"https://service.berlin.de/dienstleistung/326233/standort/327437/\",\"appointment\":{\"link\":\"\",\"slots\":\"0\",\"external\":\"false\",\"multiple\":\"0\",\"allowed\":\"false\"},\"hint\":\"\"}",
%
%// [two more entries]
%],
%"d115ServiceResponsibilityAll": false,
%"d115OnlineProcessingLink": "",
%"leikaId": 99010008012000,
%"leikaName": "Fiktionsbescheinigung Ausstellung",
%"leikaGruppe": "Aufenthaltstitel",
%"leikaKennung": "Fiktionsbescheinigung",
%"leikaVerrichtung": "Ausstellung",
%"leikaVerrichtungDetail": "",
%"ssdsName": [ "Fiktionsbescheinigung Ausstellung", "Fiktionsbescheinigung ", "Fiktionsbescheinigung"
%],
%"ssdsLemma": [
%"vorläufig", "aufenthaltstitel",
%"ausstellung", "visum", "aufenthaltserlaubnis", "fiktionsbescheinigung"
%],
%"ssdsLongName": "Ausstellung Aufenthaltstitel Fiktionsbescheinigung",
%"ssdsGruppe": [
%"Aufenthaltstitel"
%],
%"ssdsGruppeDict": [
%"Aufenthaltstitel"
%],
%"ssdsKennung": [
%"Fiktionsbescheinigung"
%],
%"ssdsKennungDict": [
%"Fiktionsbescheinigung"
%],
%"ssdsVerrichtung": [
%"Ausstellung"
%],
%"ssdsVerrichtungDict": [
%"Ausstellung"
%],
%"ssdsSynonym": [
%"Aufenthaltserlaubnis", "Aufenthaltstitel", "vorläufig", "Visum", "Fiktionsbescheinigung"],
%"ssdsSynonymDict": [
%"Aufenthaltserlaubnis", "Aufenthaltstitel", "vorläufig", "Visum", "Fiktionsbescheinigung"],
%"ssdsManualKeywords": [
%"Visa"
%],
%"_version_": 1598406083041296400
%},
%\end{minted}












%%%%%%%%%%%%%%%%%%%%%%%%%%%%%%%%%%%%%%%%%%%%%%%%%%%%%%%%%%%%%%%%%%%%%%%%%%%%%
\clearpage

%\todo{check tcolorbox and mited manuals for headings etc}
%\begin{tcblisting}{<tcb options>,
%		minted language=<language>,
%		minted style=<style>,
%		minted options={<option list>} }
%	<code>
%\end{tcblisting}




%\todo{3ashan beyetghayarou bsr3a}
% DO NOT INDENT OR PLAY AROUND HERE LEST THE CODE ZERSCHIESST SICH
\begin{figure}[H]
	\caption[HTTPS Request Body Syntax]{Request Body (26.04.2018) }
	\label{jsonFromAlexa}

\begin{minted}[linenos,tabsize=2, bgcolor=bgkolor, breaklines, fontsize=\footnotesize]{json}
{"version": "1.0",
 "session": {
  "new": true,
  "sessionId": "amzn1.echo-api.session.[unique-value]",
  "application": {
   "applicationId": "amzn1.ask.skill.[unique-value]"},
  "attributes": {
   "key": "string value"},
  "user": {
   "userId": "amzn1.ask.account.[unique-value]",
   "accessToken": "Atza|AAAAAAAA...",
   "permissions": {
   "consentToken": "ZZZZZZZ..."}}},
 "context": {
 "System": {
  "device": {
   "deviceId": "string",
   "supportedInterfaces": {
    "AudioPlayer": {}}},
  "application": {
   "applicationId": "amzn1.ask.skill.[unique-value]"},
  "user": {
   "userId": "amzn1.ask.account.[unique-value]",
   "accessToken": "Atza|AAAAAAAA...",
   "permissions": {
    "consentToken": "ZZZZZZZ..."}},
  "apiEndpoint": "https://api.amazonalexa.com",
  "apiAccessToken": "AxThk..."},
 "AudioPlayer": {
  "playerActivity": "PLAYING",
  "token": "audioplayer-token",
  "offsetInMilliseconds": 0}},
 "request": {}}
	
	\end{minted}
\end{figure}

Can be used to generate tests in Lambda.\\
source: \textsc{ASK Documentation} \footnote{\t{a\t{sk}}\href{https://developer.amazon.com/docs/custom-skills}{\lstinline|/request-and-response-json-reference.html\#request-body-syntax|}}




%###################################################################################
%###################### Appendix A          ########################################
%###################################################################################
%uncomment, if desired
\newpage
\addcontentsline{toc}{section}{Abbreviations}
\section*{List of Acronyms and Abbreviations }
\begin{flushleft}
\begin{tabular}{ll}
%\textbf{AES}	&	Advanced Encryption Standard (Symmetrisches Verschlüsselungsverfahren)\\

%\textbf{dpi}	&	dots per intch (Punkte pro Zoll; Maß für Auflösung von Bilddateien)\\
%\textbf{HTML}	&	Hypertext Markup Language (Textbasierte Webbeschreibungssprache)\\
%%\textbf{JAP}	&	Java Anon Proxy\\
%\textbf{JPEG}	&	Joint Photographic Experts Group (Graphic format)\\
%\textbf{JPG}	&	Joint Photographic Experts Group (Graphic format; short form)\\
%\textbf{LED}	&	Light Emitting Diode (lichtemittierende Diode)\\
%\textbf{LSB}	&	Least Significant Bit\\
%\textbf{MD5}	& Message Digest (Kryptographisches Fingerabdruckverfahren)\\
%\textbf{MPEG}	&	Moving Picture Experts Group (Video- einschließlich Audiokompression)\\
%\textbf{MP3}	&	MPEG-1 Audio Layer 3 (Audiokompressionformat)\\
%\textbf{PACS}	&	Picture Archiving and Communication Systems\\
%\textbf{PNG}	&	Portable Network Graphics (Grafikformat)\\
%\textbf{RSA}	&	Rivest, Shamir, Adleman (asymmetrisches Verschlüsselungsverfahren)\\
%\textbf{SHA1}	&	Security Hash Algorithm (Kryptographisches Fingerabdruckverfahren)\\
%\textbf{WAV}	&	Waveform Audio Format (Audiokompressionsformat von Microsoft)\\


\textbf{ACID}	&	Atomicity, Consistency, Isolation, Durability\\ %still needed?
\textbf{AI}		&	Artificial Intelligence\\
\textbf{ARN}	&	Amazon Resource Name\\
\textbf{ASCII}&	American Standard Code for Information Interchange\\% (Computer-Textstandard)\\
\textbf{ASK}	&	Alexa Skills Kit\\
\textbf{AVS}	&	Alexa Voice Service\\
\textbf{AWS}	&	Amazon Web Serivces\\

\textbf{B2B}	&	business-to-business\\
\textbf{B2C}	&	business-to-consumer\\

\textbf{CEO}	&	Chief Executive Officer\\
\textbf{CIA}	&	Competitive Innovation Advantage\\
\textbf{CLI}	&	Command-Line Interface\\
\textbf{CSS}	&	Cascading Style Sheets\\

\textbf{DCG}	&	Discounted Cumulative Gain\\

\textbf{FAQ}	&	Frequently Asked Question\\

\textbf{GB}		&	Gigabyte\\
\textbf{GUI}	&	Graphical User Interface\\

\textbf{HCI}	&	Human Capital Index\\
\textbf{HTML}	&	Hyper-Text Markup Language\\

\textbf{IoT}	&	Internet of Things\\
\textbf{IVR}	&	Interactive Voice Response\\

\textbf{JSON}	&	JavaScript Object Notation\\

\textbf{kbps}	&	kilobits per second\\

\textbf{LeiKa}	&	Leistungskatalog der öffentlichen Verwaltung \\
& (public administration service catalogue)\\


\textbf{ML}		&	Machine Learning\\
\textbf{MP3}	&	Moving Picture Experts Group Layer-3 Audio (MPEG-3)\\
\textbf{MVC}	&	Model-View-Controller\\
\textbf{MVP}	&	Minimum Viable Product\\

\textbf{NLP}	&	Natural Language Processing\\
\textbf{NLU}	&	Natural Language Understanding\\
\textbf{NPM}	&	Node Package Manager\\

\textbf{OS}		&	Operating System\\
\textbf{OSI}	&	Online Services Index\\
\textbf{ROI}	&	Return on Investment\\
\textbf{SSML}	&	Speech Synthesis Markup Language\\

\textbf{TF-IDF}	&	Term Frequency - Inverted Document Frequency\\
\textbf{TII}	&	Telecommunication Infrastructure Index\\

\textbf{UN}		&	United Nations\\

\textbf{VUI}	&	Voice User Interface\\


%\textbf{abk}	&	erklärung\\
%\textbf{abk}	&	erklärung\\
%\textbf{abk}	&	erklärung\\
%\textbf{abk}	&	erklärung\\
%\textbf{abk}	&	erklärung\\


\end{tabular}
\end{flushleft}


%next page
\begin{flushleft}
	\begin{tabular}{ll}
		
\textbf{appx}	&	approximately\\
\textbf{app}	&	(Mobile or web) application\\
\textbf{DevOps}	&	Development and Operations\\
\textbf{excl.}	& 	excluding\\
\textbf{Inc.}	&	incorporation\\
\end{tabular}
\end{flushleft}


%###################################################################################
%###################### Appendix B          ########################################
%###################################################################################
\newpage
%change this
\addcontentsline{toc}{section}{Glossary}
\section*{Glossary}
%%remove this
%%###################################################################################
%###################### HowTo               ########################################
%###################################################################################
\subsection*{How to Use This Template}
\begin{itemize}
\item Remove all of my text which is mostly labeled with DELETEME
\item Change the information in the 00a\_title\_page.tex file
\item Use the information written in this section
\item Ask you supervizor to help you
\item If I am not your supervizor and noone else can help you, write me an email (aubrey.schmidt@dai-labor.de)
\end{itemize}

%###################################################################################
%###################### Citations           ########################################
%###################################################################################
\subsection*{Citations}
Citing is one of the essential points you need to do in you thesis. Statements not basing on results of your own research\footnote{in what ever context} not being cited represent a breach on the rules of scientific working. Therefore, you every statement needs to be cited basing on information that other people can cross-check. A common way of citing in technical papers is: 
\begin{itemize}
\item Oberheide et al.~\cite{oberheide:2008:cloudav} state that the average time for an anti-virus enginge to be updated with a signature to detect an unknown threat is 48 days.
\end{itemize}
Note: et al. is used when the paper was written by more than two people. Check the code of this section to learn how to cite from a technical perspective.

Note: you can change the citation style in the \texttt{thesis.tex} file, e.g. to harvard style citations. Instructions on this can also be found in this file.

You should not cite anything that can be changed, e.g. it is not that good citing web pages since they might get updated changing the cited content. There are no clear quality measures on citing sources but aubrey believes that the following list is true for several cases, starting with highest quality:
\begin{enumerate}
\item Journal article or book
\item Conference paper
\item Workshop paper
\item Technical report
\item Master thesis
\item Bachelor thesis
\item General Web reference
\end{enumerate}
There might be workshop papers that have a higher quality than some journal papers. Therefore this list only gives you a hint on possible quality measures. Another measure can be whether a paper was indexed by ACM/IEEE, although this is not a strong indicator.

%###################################################################################
%###################### Papers             ########################################
%###################################################################################
\subsection*{Finding and Handling Citation Sources}
Following ressources are required for finding and handling articles, books, papers and sources.
\begin{itemize}
\item your primary resource will be \url{http://scholar.google.com}
\item \url{http://www.google.com} might also be used
\item \url{wikipedia.com} can be a good start for finding relevant papers on your topic
\item you should download and install JabRef or a similar tool \url{http://jabref.sourceforge.net/}
\item you should point JabRef to the mybib.bib file
\item you should immediately enter a relevant paper to JabRef, additionally, you should write a short summary on it; else, you will do this work at least twice.
\end{itemize}

%###################################################################################
%###################### General             ########################################
%###################################################################################
\subsection*{General Advices}
\begin{itemize}
\item Do not take care of design, \LaTeX will do this for you. If you still feel that you need to take care of this, do this when you have finished writing, else you will end up in a lot of double and triple work.
\item \LaTeX will do exactly that you will tell it to do. If you have problems with this, go for google or ask you supervizor
\item use labels in order to be able to reference to chapters, section, subsections, figures, tables, etc. ...
\end{itemize}

%###################################################################################
%###################### Commands            ########################################
%###################################################################################
\subsection*{General Commands}
\begin{itemize}
\item check \url{http://en.wikibooks.org/wiki/LaTeX}
\item check \url{http://www.uni-giessen.de/hrz/tex/cookbook/cookbook.html} German
\end{itemize}
Please also check the following source~\cite{latexcookbook2007}.

%###################################################################################
%###################### Code                ########################################
%###################################################################################
\newpage
\subsection*{Code}
This section shows you how to get your code into a \LaTeX document. See code for options.
\lstinputlisting[language=JAVA,xleftmargin=8mm,]{__help/Example.java}


\lstset{ %
language=Java,   	             % the language of the code
numbers=left,
basicstyle=\scriptsize,       % the size of the fonts that are used for the code
numbers=left,                   % where to put the line-numbers
numberstyle=\scriptsize,      % the size of the fonts that are used for the line-numbers
stepnumber=1,                   % the step between two line-numbers. If it's 1, each line 
                                % will be numbered
numbersep=5pt,                  % how far the line-numbers are from the code
backgroundcolor=\color{white},  % choose the background color. You must add \usepackage{color}
showspaces=false,               % show spaces adding particular underscores
showstringspaces=false,         % underline spaces within strings
showtabs=false,                 % show tabs within strings adding particular underscores
frame=single,                   % adds a frame around the code
tabsize=2,                      % sets default tabsize to 2 spaces
captionpos=b,                   % sets the caption-position to bottom
breaklines=true,                % sets automatic line breaking
breakatwhitespace=false,        % sets if automatic breaks should only happen at whitespace
title=\lstname,                 % show the filename of files included with \lstinputlisting;
                                % also try caption instead of title
xleftmargin=8mm,
framexleftmargin=4mm,           
escapeinside={\%*}{*)},         % if you want to add a comment within your code
morekeywords={*,...}            % if you want to add more keywords to the set
}
\begin{lstlisting}[float=h, caption=Example code is presented here, label=list:code, frame=single]
class Beispiel{

	public static void main(String args[]){
	
		System.out.println("Hello World");
		
	}
	
}
\end{lstlisting}

%###################################################################################
%###################### Figures             ########################################
%###################################################################################
\newpage
\subsection*{Figures}
This section describes how to include images to your document. Information was taken from \url{http://en.wikibooks.org/wiki/LaTeX/Floats,_Figures_and_Captions}, visited on 05/08/2011. Please make sure to use original vector graphics as basis since image quality might be used as weak indicator for thesis quality. For this, try to find find files in \texttt{.SVG} or \texttt{.PDF} format. Exporting a \texttt{.PNG} or \texttt{.JPG} to \texttt{.PDF} will not work since data was already lost while exporting it to these formats. This is the case for most Web graphics. Wikipedia startet entering most in images in \texttt{.SVG} which easily can be transformed to \texttt{.PDF}, but please do not forget proper citations.

%use a modifier to decide where you desire Latex should put the image: (h)ere, (t)op, (b)ottom, or nothing more than that image on a (p)age
\begin{figure}[h]%[htbp]
%this will center your image
\centering
%this will include your image
\includegraphics[width=0.15\textwidth]{template/TUBerlin_Logo_rot_hell}
%this is the caption + label. the label will not be printed in the caption. Moving the label out of the caption can result in problems.
\caption[Including an Image]{Including an image; in this case a PDF. Please note that the caption is placed below the image.\label{fig:help1}}
\end{figure}

\begin{figure}[h]
%this will center your image
\centering
%this will include your image
\includegraphics[width=0.25\textwidth]{template/aot_logo}
%this is the caption + label. the label will not be printed in the caption. Moving the label out of the caption can result in problems.
\caption[Short caption for list of figures]{See code for caption options: this is a long caption which is printed in the Text. Additionally, image size was increased\label{fig:help2}}
\end{figure}


\begin{figure}[h]
  \centering
  \subfloat[Small]{\label{fig:tub1}\includegraphics[width=0.1\textwidth]{template/TUBerlin_Logo_rot_hell}}                
  \subfloat[Large]{\label{fig:tub2}\includegraphics[width=0.3\textwidth]{template/TUBerlin_Logo_rot_hell}}
  \subfloat[Medium]{\label{fig:tub3}\hspace{2cm}\includegraphics[width=0.2\textwidth]{template/TUBerlin_Logo_rot_hell}}
  \caption[Placing images side by side]{Placing images side by side using the subfig package. Space between the images can be adjusted.\label{fig:tuball}}
\end{figure}


%###################################################################################
%###################### Tables              ########################################
%###################################################################################
\newpage
\subsection*{Tables}
Here, you will find some example tables.The tables were taken from \url{http://en.wikibooks.org/wiki/LaTeX/Tables}, visited on 05/08/2011. Table environment was added plus caption and label. For code, check \url{__help/latex_hinweise.tex}.

\begin{table}[h]
\caption{Simple table using vertical lines. Note that the caption is always above the table! Please check code for finding the right place for the table label.\label{tab:help1}}
\centering
  \begin{tabular}{ l | c || r | }
    \hline
    1 & 2 & 3 \\ \hline
    4 & 5 & 6 \\ \hline
    7 & 8 & 9 \\
    \hline
  \end{tabular}
\end{table}

\begin{table}[h]
\caption{Table using vertical and horizontal lines\label{tab:help2}}
\centering
\begin{tabular}{|r|l|}
  \hline
  7C0 & hexadecimal \\
  3700 & octal \\ \cline{2-2}
  11111000000 & binary \\
  \hline \hline
  1984 & decimal \\
  \hline
\end{tabular}
\end{table}

\begin{table}[h]
\caption{Table with column width specification on last column\label{tab:help3}}
\centering
    \begin{tabular}{ | l | l | l | p{5cm} |}
    \hline
    Day & Min Temp & Max Temp & Summary \\ \hline
    Monday & 11C & 22C & A clear day with lots of sunshine.  
    However, the strong breeze will bring down the temperatures. \\ \hline
    Tuesday & 9C & 19C & Cloudy with rain, across many northern regions. Clear spells
    across most of Scotland and Northern Ireland,
    but rain reaching the far northwest. \\
    \hline
    \end{tabular}
\end{table}

\begin{table}[h]
\caption{Table using multi-column and multirow\label{tab:help4}}
\centering
\begin{tabular}{|l|l|l|}
\hline
\multicolumn{3}{|c|}{Team sheet} \\
\hline
Goalkeeper & GK & Paul Robinson \\ \hline
\multirow{4}{*}{Defenders} & LB & Lucus Radebe \\
 & DC & Michael Duberry \\
 & DC & Dominic Matteo \\
 & RB & Didier Domi \\ \hline
\multirow{3}{*}{Midfielders} & MC & David Batty \\
 & MC & Eirik Bakke \\
 & MC & Jody Morris \\ \hline
Forward & FW & Jamie McMaster \\ \hline
\multirow{2}{*}{Strikers} & ST & Alan Smith \\
 & ST & Mark Viduka \\
\hline
\end{tabular}
\end{table}






\begin{flushleft}
\begin{tabular}{ll}

\textbf{Access Key}			&	 	\ref{accesskeys}\\


\textbf{Alexa}				&		\ref{alexa:def}\\
\textbf{Alexa Skill}		&		\ref{alexa:def}\\
\textbf{Alexa Skills Kit}	&		\ref{ask:def}\\
\textbf{Amazon Developer Console}&	\ref{ask:devconsole}\\

\textbf{Application ID}		&	required by Lambda. Skill ID in our context\\

\textbf{Intent}				&		\ref{intents}\\
\textbf{Interaction Model}	&		\ref{interactionMdl:def}, \ref{interactionModel}\\



\textbf{Lambda, AWS Lambda Function}	&		\ref{aws:modules}\\
%\textbf{Lambda Function}	&		erklärung\\
\textbf{Lex, Amazon Lex}	&		\ref{aws:modules}\\

\textbf{Node.js}			&		\ref{nodejs:def}\\

\textbf{Polly, Amazon Polly}&		\ref{aws:modules}\\
\textbf{Public Service}			&	\ref{pubsvc}\\
\textbf{Slot}				&		\ref{slots}\\
\textbf{Skill ID}			&	amzn1.ask.skill.d7732837-fab2-42ff-a152-4eb0fc4ee646\\
\textbf{Transcribe, Amazon Transcribe}	&		\ref{aws:modules}\\

\textbf{Utterance}			&		\ref{utterances}\\

%\textbf{ElasticSearch}		&		erklärung\\


\textbf{Web Service}			&		e.g. AWS\\




\textbf{Bot}				&	Unless otherwise mentioned, a chatbot\\
\textbf{Hitlist}			&	erklärung\\
\textbf{Voice-first}		& \shortstack[l]{an always-on, intelligent piece of hardware, where the \\ primary interface is voice, both input and output.}\\

%\textbf{abk}				&	erklärung\\
%\textbf{abk}				&	erklärung\\
%\textbf{abk}				&	erklärung\\
%\textbf{abk}				&	erklärung\\
%\textbf{abk}				&	erklärung\\


\end{tabular}
\end{flushleft}
